\section{Cơ sở lý thuyết và công trình liên quan}
\label{sec:related}

\subsection{Leo thang đặc quyền trong multi-tenant Kubernetes}
Trong cụm Kubernetes nhiều người thuê, các tenant được cô lập mức logic bằng
namespace nhưng chia sẻ chung nhân của nút. Bề mặt nguy hiểm nhất là \emph{leo
thang đặc quyền} và \emph{thoát container}: tận dụng lỗ hổng nhân hoặc cấu hình
sai (capabilities dư thừa, mount host-path, \texttt{ptrace}, \texttt{unshare}/
\texttt{setns}, nạp mô-đun nhân) để vượt ranh giới tenant và chiếm quyền trên
nút~\cite{ref5_combe_docker}. Mọi con đường này đều biểu hiện thành các
\emph{chuỗi syscall đặc quyền} đặc trưng, nên phân tích syscall là điểm quan sát
tự nhiên để phát hiện chúng.

\subsection{Bảo mật lời gọi hệ thống cho container}
Trong mô hình chia sẻ nhân, nhân Linux phơi bày hàng trăm syscall, phần lớn không
cần thiết cho một khối lượng công việc cụ thể, làm phình to bề mặt tấn
công~\cite{ref11_shu_security}. Các cơ chế truyền thống (seccomp, AppArmor,
SELinux) lọc theo chính sách tĩnh nên: (1) không thích nghi với hành vi thời gian
chạy; (2) thiếu nhận thức về \emph{chuỗi} syscall nên mù với tấn công nhiều bước;
và (3) tốn chi phí cấu hình thủ công ở quy mô lớn.

\subsection{Sinh danh sách syscall và lọc động}
Confine~\cite{ref31_confine} dùng phân tích tĩnh để rút gọn hồ sơ seccomp (SRR
$\approx 47{,}6\%$) nhưng không thích nghi thời gian chạy.
Prof-gen~\cite{ref36_profgen} kết hợp theo dõi lúc khởi động với đồ thị lời gọi
tĩnh, vẫn bỏ sót thay đổi hành vi về sau. Optimus~\cite{ref27_optimus} và
KubeRosy~\cite{ref29_kuberosy} dùng eBPF để cập nhật chính sách thời gian chạy,
song Optimus cần phê duyệt thủ công và KubeRosy gắn chặt với Kubernetes.
NIMOS~\cite{ref38_nimos} đánh giá giá trị của lọc dựa trên chuỗi. DeSFAM định vị
mình ở chỗ hợp nhất cả ba: lập hồ sơ lai, phát hiện bất thường học máy, và thực
thi trong nhân.

\subsection{Phát hiện xâm nhập dựa trên chuỗi syscall}
Hướng phát hiện bất thường dựa trên chuỗi syscall có lịch sử lâu dài, khởi nguồn
từ ``a sense of self''~\cite{forrest1996sense,hofmeyr1998ids} và các mô hình thay
thế~\cite{warrender1999detecting}, cùng cảnh báo kinh điển về tấn công bắt
chước~\cite{wagner2002mimicry}. Các hướng hiện đại gồm phương pháp ngữ
nghĩa~\cite{creech2014semantic}, túi lời gọi (BoSC) cho
container~\cite{abed2015bag}, và học sâu một lớp/biến phân như
DAGMM~\cite{zong2018dagmm}, Deep SVDD~\cite{ruff2018deepsvdd},
HIDS học sâu~\cite{haider2021deep}. Riêng cho container có các đánh giá phát hiện
bất thường~\cite{castanhel2021peek} và ngữ cảnh hóa syscall~\cite{elkhairi2022chids},
cũng như tập hợp Random/Isolation Forest~\cite{bsoul2022ensemble}. SyscallAD của
DeSFAM nằm trong dòng này: hợp nhất VAE~\cite{vae_kingma} (lỗi tái cấu trúc) với
iForest~\cite{iforest_liu} (ngoại lệ thống kê), bổ sung khớp mẫu lành tính bằng
PrefixSpan~\cite{prefixspan_pei}.

\subsection{Bộ dữ liệu DongTing}
DongTing~\cite{ref26_dongtingdataset} là bộ dữ liệu quy mô lớn gồm các chuỗi
syscall của nhân Linux được gán nhãn bình thường/độc hại (theo các bản vá lỗi
syzkaller), thích hợp cho huấn luyện và xác nhận bộ phát hiện bất thường. Chuyên
đề này dùng DongTing làm nguồn dữ liệu ngoại tuyến duy nhất, đúng như DeSFAM.

\subsection{eBPF và Tetragon}
eBPF cho phép nạp logic giám sát/thực thi an toàn vào nhân tại các điểm vào syscall
với chi phí thấp~\cite{ref20_bpf_security}. DeSFAM dùng eBPF/LSM để chặn nội tuyến.
Trong chuyên đề này, thay vì tự viết chương trình eBPF, chúng tôi dùng
Tetragon~\cite{tetragon} (thuộc hệ sinh thái Cilium~\cite{cilium})---một lớp quan
sát và thực thi bảo mật dựa trên eBPF cho Kubernetes---làm \emph{cảm biến syscall}
phát sự kiện qua gRPC. Đây là lựa chọn thực dụng để hiện thực hóa lớp thu thập của
DeSFAM mà không cần sửa nhân.

\subsection{Vị trí của chuyên đề}
Chuyên đề kế thừa DeSFAM như một \emph{công trình tham chiếu} cho thiết kế lõi phát
hiện bất thường syscall, và vận dụng nó để xây dựng một \emph{cơ chế phát hiện leo
thang đặc quyền} cho ngữ cảnh multi-tenant Kubernetes---một ngữ cảnh mà DeSFAM
không đặt trọng tâm. Đóng góp nằm ở chỗ: (i) định vị các đặc trưng syscall đặc
quyền làm tín hiệu cho leo thang đặc quyền; (ii) hiện thực hóa cơ chế thành một bộ
phát hiện trực tiếp giám sát theo namespace trên Tetragon/Kubernetes, lấp khoảng
trống ``từ mô phỏng đến triển khai''; và (iii) đối chiếu trung thực hiệu năng với
số liệu công bố của công trình tham chiếu.
